\documentclass[../main.tex]{subfiles}

\begin{document}

% ===========================================================================
\subsection{Conclusiones}
% ===========================================================================

Este proyecto partió de una pregunta aparentemente sencilla: ¿puede una base de datos médica ser, al mismo tiempo, conforme con el RGPD por diseño y suficientemente eficiente para uso clínico en tiempo real? El trabajo presentado en la Parte~I ofrece una respuesta concreta y empíricamente validada: sí es posible, siempre que las decisiones arquitectónicas correctas se tomen a nivel del esquema de datos y no se deleguen en código de aplicación o en capas de control de acceso. A partir de ahí, el proyecto creció hasta abarcar tres paradigmas complementarios ---relacional, semiestructurado y documental---, lo que permitió comprobar de primera mano que no hay una única tecnología \enquote{correcta}, sino la más adecuada para cada tipo de dato clínico.

\paragraph{Sobre el cumplimiento de la privacidad por diseño.}
La separación basada en SHA-256 entre \texttt{paciente\_anonimo} y \texttt{datos\_personales} alcanza una seudonimización estructural genuina en el sentido del Considerando~26 del RGPD. La propiedad crítica es que la anonimización no es una capa extraíble: al usuario \texttt{investigador} se le deniega el acceso a \texttt{datos\_personales} a nivel del motor de base de datos, no a través de lógica de aplicación que pudiera ser eludida. Relacionar un \texttt{codigo\_hash} con una identidad real requiere o bien el NIF original (que reside fuera de la base de datos) o un ataque de preimagen contra SHA-256, actualmente considerado computacionalmente inviable. Las restricciones \texttt{ON DELETE RESTRICT} garantizan adicionalmente que ningún registro pueda eliminarse del modelo clínico mientras exista su contrapartida de datos personales, preservando la integridad del historial de auditoría.

\paragraph{Sobre la optimización multi-motor.}
Los tres motores de almacenamiento elegidos para ChronoHealthDB ofrecieron beneficios medibles en sus respectivos roles:

\begin{itemize}
    \item \textbf{InnoDB} proporcionó las garantías transaccionales, la integridad referencial y el control de concurrencia multiversión necesarios para las tablas del flujo de trabajo clínico. El benchmarking confirmó que los índices compuestos (\texttt{idx\_paciente\_fecha}) redujeron los tiempos de join del historial de paciente de 39,89~ms a 0,64~ms (\textbf{mejora de 62,3×}), y que el índice temporal (\texttt{idx\_fecha\_evento}) redujo las consultas de conteo mensual de 46,06~ms a 1,36~ms (\textbf{33,8×}).

    \item \textbf{ARCHIVE} demostró que la optimización del almacenamiento para registros históricos inmutables no es una preocupación marginal: el mismo conjunto de datos que ocupa 2.560~KiB bajo InnoDB se comprime a 533~KiB bajo ARCHIVE, una \textbf{reducción del 79,2\%} atribuible a la compresión \textit{zlib} fila a fila y a la ausencia de la infraestructura de gestión MVCC. A escala de una consejería de salud regional que gestione 200.000 bajas médicas al año durante una década, esto equivale a más de 640~GB de diferencia en almacenamiento primario.

    \item \textbf{MEMORY} confirmó su idoneidad para la gestión de colas en tiempo real: las consultas de prioridad de triaje se ejecutaron en 0,77~ms, un \textbf{45,9\% más rápido} que la tabla InnoDB equivalente. Más importante aún, MEMORY garantiza una latencia constante independiente de la carga en disco, propiedad que InnoDB no puede ofrecer cuando su \textit{buffer pool} está bajo presión.
\end{itemize}

\paragraph{Sobre la generación de datos sintéticos.}
El conjunto de datos sintético de 40.000 pacientes, con plantillas de texto orientadas a escenarios clínicos, constantes vitales correlacionadas con la patología y factores de riesgo estratificados por edad, validó que los resultados del benchmarking no son artefactos de datos trivialmente simples. El conjunto de datos se comporta como datos reales desde la perspectiva de los patrones de consulta, lo que refuerza la credibilidad de las mediciones de rendimiento.

\paragraph{Sobre la fase semiestructurada (Parte~II).}
Llevar los mismos datos a XML y cargarlos en eXist-db demostró que la naturaleza jerárquica del historial clínico se aprovecha mejor cuando se almacena como un árbol que cuando se reparte en tablas. Las tres consultas XQuery funcionaron como un sistema de soporte a la decisión real: la triangulación cruzó urgencias, historial y bajas en una sola consulta FLWOR; la detección de crónicos se reescribió para pasar de un coste cercano al cuadrático a uno lineal apoyándose en la propia estructura del documento; y el monitor de hipoxia obligó a manejar con cuidado las conversiones de tipos (\texttt{castable as xs:double}), un detalle nada trivial al trabajar con datos que en XML son siempre texto. La lección es que el esfuerzo invertido en dar buena forma a los documentos se recupera con creces a la hora de consultarlos.

\paragraph{Sobre la arquitectura híbrida con MongoDB (Parte~III).}
La fase NoSQL confirmó que MySQL y MongoDB pueden convivir asignando a cada uno lo que mejor hace: lo transaccional al relacional y el conocimiento clínico flexible (las guías) al documental. Tres ideas quedaron especialmente claras. Primera, que la privacidad desde el diseño es viable de verdad: al anonimizar el DNI con SHA-256 en la propia inserción, los 50.000 expedientes llegaron a MongoDB ya seudonimizados, sin ventana alguna de exposición. Segunda, que el análisis con \texttt{.explain()} reproduce en un motor NoSQL la misma lección de los índices que vimos en MySQL: el índice \texttt{anio\_nacimiento\_1} transformó un \texttt{COLLSCAN} de 50.000 documentos en un \texttt{IXSCAN} que solo examina los 12.040 relevantes. Y tercera, que delegar el procesamiento masivo en el motor (los \textit{aggregation pipelines}) y reservar Python como capa de presentación es la decisión arquitectónica que separa una solución que escala de un anti-patrón que satura la memoria del cliente.

\paragraph{Valoración global.}
El proyecto demuestra que la tensión entre el cumplimiento del RGPD y el rendimiento operativo es un falso dilema. El cumplimiento por diseño, mediante separación estructural de datos, y la optimización del rendimiento, mediante selección de motor e indexación, no son objetivos contrapuestos sino decisiones arquitectónicas complementarias que se refuerzan mutuamente. Más allá de lo relacional, recorrer los tres paradigmas sirvió para entender que la elección de la tecnología de almacenamiento no es una cuestión de modas, sino de ajustar la herramienta a la forma del dato: tablas para lo tabular, árboles XML para lo jerárquico y documentos BSON para lo flexible. La complejidad añadida de manejar tres tecnologías queda plenamente justificada por lo que cada una aporta donde las demás flaquean.

% ===========================================================================
\subsection{Líneas Futuras de Investigación}
% ===========================================================================

La fase relacional presentada en este trabajo constituye la primera capa de un sistema naturalmente extensible. Varias líneas de investigación emergen directamente de sus limitaciones actuales:

\paragraph{Bases de datos de grafos para las relaciones entre pacientes.}
Las Partes~II y~III ya cubrieron los paradigmas semiestructurado (XML) y documental (MongoDB), de modo que el siguiente salto natural es el de las \textbf{bases de datos de grafos} (Neo4j o Apache AGE sobre PostgreSQL). Las redes clínicas son inherentemente relacionales en el sentido de la teoría de grafos: redes de comorbilidades, grafos de interacciones entre fármacos, análisis de vías de atención o trazado de contactos epidémicos son consultas que se expresan de forma natural como recorridos de grafo pero resultan costosas como operaciones SQL con múltiples joins o incluso como agregaciones documentales. Una extensión interesante consistiría en proyectar la tabla \texttt{evento\_clinico} sobre un grafo de propiedades y evaluar algoritmos como el camino más corto o la detección de comunidades sobre los datos ya seudonimizados.

\paragraph{Texto clínico libre y búsqueda semántica.}
El modelo actual estructura bien los eventos, pero las notas clínicas de texto libre, los informes de imagen y los informes de alta siguen siendo el gran reto. Una línea natural sería explotar las capacidades de búsqueda de texto completo de MongoDB ---o incorporar índices vectoriales y modelos de lenguaje--- para hacer búsqueda semántica sobre ese contenido no estructurado, manteniendo siempre la seudonimización por hash que ya garantiza el resto del sistema.

\paragraph{Escalabilidad horizontal y replicación.}
El sistema actual es una instancia MySQL de nodo único. Un despliegue en producción a escala hospitalaria requeriría una topología de replicación (InnoDB Cluster o MySQL Group Replication) con réplicas de lectura para la carga de trabajo de investigación y un nodo primario para las escrituras clínicas. El efecto del retraso de replicación sobre las tablas MEMORY —que son locales a cada nodo— es una cuestión de ingeniería abierta que merece investigación específica.

\paragraph{Cifrado de tablas ARCHIVE.}
El cifrado de tablespace de InnoDB (\texttt{ENCRYPTION='Y'}) está bien documentado, pero el cifrado del motor ARCHIVE no tiene soporte nativo en todas las configuraciones. Explorar el cifrado transparente a nivel de disco para tablas ARCHIVE, o migrar la capa histórica a un motor compatible con cifrado y compresión simultáneos, abordaría el requisito regulatorio de protección de datos en reposo sin sacrificar los beneficios de compresión demostrados en este trabajo.

\paragraph{Aprendizaje automático sobre datos anonimizados.}
El conjunto de 40.000 pacientes contiene ya señal suficiente para entrenar modelos predictivos básicos sobre datos seudonimizados: riesgo de reingreso, estimación de duración del ingreso o puntuaciones de alerta temprana basadas en las constantes vitales del triaje. Una \textit{pipeline} de aprendizaje automático que opere exclusivamente sobre las tablas accesibles al usuario \texttt{investigador} demostraría una aplicación real y concreta de la arquitectura conforme con el RGPD diseñada en este trabajo.

\end{document}
